% =============================================================================
% Multimodal deepfake detection: IEEE conference-format paper, FIRST VERSION.
% Class: IEEEtran, [conference] option (two-column, 10 pt). Overleaf compatible:
% no absolute paths, only packages shipped with TeX Live / MiKTeX.
% Build: latexmk -pdf main.tex   (BibTeX with IEEEtran.bst, refs.bib)
%
% STATE OF THIS FILE (2026-10-24): STRUCTURE ONLY. No experiment has been run.
% Rules in force (CLAUDE.md):
%   - No number appears here unless it comes from a logged experiment E-xxx.
%     Every sentence that will carry a result must keep a marker "% E-xxx".
%   - Cited numbers from other papers are named with their own protocol and
%     are never presented as ours.
%   - Only papers in docs/RESEARCH.md with status READ are cited.
%   - No "real-time", "state of the art", "better than X" without a
%     head-to-head run on the same protocol.
%   - No em-dashes or en-dashes (do not write double or triple hyphens).
%   - Model, framework, dataset, metric choices are OPEN (docs/DECISIONS.md,
%     O-001 to O-011) and appear below only as \tbd{...} placeholders.
% =============================================================================
\documentclass[conference]{IEEEtran}

\usepackage[T1]{fontenc}
\usepackage[utf8]{inputenc}
\usepackage{amsmath,amssymb}
\usepackage{booktabs}
\usepackage{graphicx}
\usepackage{xcolor}
\usepackage{cite}
\usepackage[hidelinks]{hyperref}

% Visible placeholder for anything that is not yet decided or measured.
\newcommand{\tbd}[1]{\textbf{\textcolor{red}{[TBD: #1]}}}
% Visible placeholder for a number that must come from a logged experiment.
\newcommand{\fromexp}[1]{\textbf{\textcolor{red}{[to be filled from #1]}}}

\begin{document}

%%%%%%%%%%%%%%%%%%%%%%%%%%%%%%%%%%%%%%%%%%%%%%%%%%%%%%%%%%%%
% PROPOSED TITLES (proposals only, none is final; the authors and guide choose).
%   Proposal 1: Multimodal Deepfake Detection for Images, Video and Audio:
%               A Hypothesis-Driven Study Under a Source-Disjoint Protocol
%   Proposal 2: From a Minimal Baseline to a Measured Detector: Seven
%               Documented Iterations of Image, Video and Audio Deepfake Detection
%   Proposal 3: How Far Does It Generalise? A Multi-Seed Study of Multimodal
%               Deepfake Detection With Paired Comparisons
% Proposal 1 is used below as a working title.
%%%%%%%%%%%%%%%%%%%%%%%%%%%%%%%%%%%%%%%%%%%%%%%%%%%%%%%%%%%%
\title{Multimodal Deepfake Detection for Images, Video and Audio:\\
A Hypothesis-Driven Study Under a Source-Disjoint Protocol}

\author{%
\IEEEauthorblockN{Vaibhav Shiroorkar, Shivansh Gupta, Adithya Jithesh}
\IEEEauthorblockA{Department of Electronics and Computer Engineering\\
K. J. Somaiya School of Engineering, Somaiya Vidyavihar University\\
Vidyavihar, Mumbai 400 077, India\\
\tbd{contact e-mail addresses, not supplied}}
\and
\IEEEauthorblockN{Dr. Tejashree Phatak (Guide)}
\IEEEauthorblockA{Assistant Professor\\
Department of Electronics and Computer Engineering\\
K. J. Somaiya School of Engineering, Somaiya Vidyavihar University\\
Vidyavihar, Mumbai 400 077, India}}

\maketitle

%%%%%%%%%%%%%%%%%%%%%%%%%%%%%%%%%%%%%%%%%%%%%%%%%%%%%%%%%%%%
\begin{abstract}
% Abstract skeleton. Each bracket is filled only after the experiment exists.
Synthetic or manipulated faces, videos and voices can be produced with
consumer hardware, and published detectors are usually compared under
different protocols, which makes their reported numbers hard to read side by
side. This paper describes a detector for images, videos and the audio of
videos that is built from a minimal baseline through a fixed number of
iterations, each started from a written hypothesis and a stated failure
criterion. All candidates are run in one pipeline on a split that keeps
identities and source recordings apart, with several random seeds and
interval estimates that resample independent units.
\tbd{one sentence on the final system, after the model and dataset choices
O-003, O-006, O-007 are made}.
Results: \fromexp{E-xxx, final comparison}.
\tbd{one sentence on what did not work, after the iterations are run}.
\tbd{one sentence on the limits of the evidence}.
\end{abstract}

\begin{IEEEkeywords}
deepfake detection, multimodal, generalisation, evaluation protocol,
calibration
\end{IEEEkeywords}

%%%%%%%%%%%%%%%%%%%%%%%%%%%%%%%%%%%%%%%%%%%%%%%%%%%%%%%%%%%%
\section{Introduction}\label{sec:intro}

Face swaps, talking-face synthesis and cloned voices can now be produced
with consumer hardware and off-the-shelf software~\cite{dolhansky2020dfdc}. A detector that is to be useful must handle images, video, and
the audio track of a video, and must say how sure it is. The literature
supplies many detectors and datasets, but the numbers reported for them are
measured under different protocols. In one benchmark that trained all
detectors on one dataset and tested them within and across datasets, a plain
Xception network reached a mean AUC of 0.9450 on the in-domain tests and
0.7718 on the cross-domain tests~\cite{yan2023deepfakebench}. % source: R-017, their protocol, not ours
The gap between in-dataset and cross-dataset results recurs across face
datasets~\cite{li2020celebdf,li2020facexray,dolhansky2020dfdc}, so a single
in-dataset number says little about behaviour on new material.

This work treats the path to a result as part of the result. We build the
simplest end-to-end detector first, then make a fixed number of
iterations, each with a hypothesis, a success criterion and a failure
criterion written before the run. Failed iterations are kept in the record.
The aim is not the best model in the world, but results that can be compared
with published detectors on a protocol that we define and run ourselves.

\paragraph*{Contributions}
% To be rewritten to match what the logged experiments actually support.
The intended contributions are:
\begin{enumerate}
  \item a written evaluation protocol with source-disjoint splits, a
        threshold chosen on validation data only, and at least three seeds
        per candidate (Section~\ref{sec:exp});
  \item a minimal baseline and seven documented iterations with their
        hypotheses and outcomes, including failures
        (Section~\ref{sec:results}, \fromexp{E-xxx, I-0 to I-7});
  \item a paired comparison of candidate models in one pipeline, with cost
        measured next to accuracy (\fromexp{E-xxx, comparison});
  \item a test of whether remote photoplethysmography (heart-rate) evidence
        adds anything beyond the interval of the paired difference
        (\fromexp{E-xxx, rPPG feasibility}).
\end{enumerate}

The rest of the paper is organised as follows. Section~\ref{sec:related}
reviews related work, Section~\ref{sec:method} describes the method,
Section~\ref{sec:exp} the protocol, Section~\ref{sec:results} the results,
and Sections~\ref{sec:disc} to~\ref{sec:concl} discuss, limit and conclude.

%%%%%%%%%%%%%%%%%%%%%%%%%%%%%%%%%%%%%%%%%%%%%%%%%%%%%%%%%%%%
\section{Related Work}\label{sec:related}

% Written from papers in docs/RESEARCH.md marked READ (or READ in part).
% Paraphrased. Numbers from other papers are not reproduced unless the
% protocol is named in the same sentence.

\subsection{Datasets for face forgery}
FaceForensics++ provides real videos manipulated with four methods at
several compression levels, together with a fixed split and a hidden-test
benchmark, and it reports that detectors lose accuracy as compression
increases~\cite{rossler2019faceforensics}. Celeb-DF was built with an
improved synthesis pipeline to remove visible artefacts, and the authors
evaluated existing detectors, run with their released models, on it and on
earlier sets, finding it the hardest of those compared~\cite{li2020celebdf}.
The DFDC dataset contributes a large collection of videos from consenting
actors, with held-out subjects and a test set that includes content from
outside the dataset; its analysis of the challenge entries reports weak
performance on the test content that came from outside the dataset~\cite{dolhansky2020dfdc}.

\subsection{Visual detectors and generalisation}
Face X-ray predicts the blending boundary between a pasted face and its
background and is trained on self-made blended images, so that no fakes from
a particular generator are required~\cite{li2020facexray}. DeepfakeBench
puts many detectors in one codebase with one preprocessing and metric
pipeline, and finds that plain strong backbones come close to more elaborate
methods when the training recipe is shared~\cite{yan2023deepfakebench}.
For still images in general, replacing a trained real-versus-fake network by
a linear or nearest-neighbour classifier on the frozen features of a
pretrained vision-language model generalised across generators better than
the trained baselines in the authors' own tests~\cite{ojha2024universal}.
Work on transfer learning for image classification relates accuracy on a
source task to accuracy after transfer, with caveats that depend on the
regularisers used~\cite{kornblith2019better}; whether this carries over to
forensic detection is an open question that we test.

\subsection{Physiological signals}
FakeCatcher derives photoplethysmography-like signals from facial regions
and classifies them, and reports high in-dataset accuracy on the generators
of its time~\cite{ciftci2020fakecatcher}. A follow-up uses the same kind of
signal to identify which generator produced a fake~\cite{ciftci2020hearts}.
DeepFakesON-Phys fine-tunes a heart-rate estimation network as a detector,
with one model per database and identity-disjoint development and test
sets~\cite{hernandez2020deepfakeson}. Later work is more cautious. A
study with its own recordings reports that a heart-rate signal can be
recovered from high-quality face swaps and can follow the driving
video~\cite{seibold2025heart}. An independent reproduction of
DeepFakesON-Phys on talking-face fakes, under an identity-disjoint split,
reported a much lower AUC than the original paper's in-dataset
figure~\cite{harraq2026physio}. Pulse signals can also be edited on
purpose, which makes a pulse-based check open to attack~\cite{lin2023ppg}.
These results are why we treat remote photoplethysmography as a hypothesis
to test, not as a component to adopt.

\subsection{Audio and audio-visual detection}
ASVspoof~2019 is the standard public database of synthesised, converted and
replayed speech, with attacks in the evaluation set that are absent from
training~\cite{wang2020asvspoof}. AASIST is a compact model on the raw
waveform that applies graph attention over spectral and temporal nodes, and
its authors report variation across seeds~\cite{jung2021aasist}.
FakeAVCeleb provides videos with both visual and cloned-voice manipulations,
and its baselines show that naive combinations of unimodal detectors can
perform poorly~\cite{khalid2022fakeavceleb}. A two-stream audio-visual
network that scores the dissonance between the streams over short chunks can
also localise manipulated segments~\cite{chugh2021notmade}.

\subsection{Calibration}
Modern neural networks are often poorly calibrated, and temperature scaling,
a single parameter fitted on validation data, is an effective
correction on the tasks the authors studied~\cite{guo2017calibration}. We
use their definitions of expected calibration error and temperature scaling,
and fit any such correction on validation data only.

%%%%%%%%%%%%%%%%%%%%%%%%%%%%%%%%%%%%%%%%%%%%%%%%%%%%%%%%%%%%
\section{Method}\label{sec:method}
% SKELETON. Nothing in this section is decided yet. Each placeholder names
% the open decision (docs/DECISIONS.md) that fills it.

\subsection{Problem statement}
Let $x$ denote an input (an image, a video clip, or the audio of a clip) and
$y\in\{0,1\}$ its label, with $y=1$ for manipulated content. A detector
outputs a score $s(x)\in[0,1]$. The exact output and the labels used are
\tbd{O-001, output definition}.
How image, video and audio inputs are each handled is \tbd{O-002}.

\subsection{Data and splits}
Datasets: \tbd{O-003}. Each dataset's licence, source and hash are recorded
before use. Every split is stored as a file with an identifier and a hash,
and no identity, source video or near-duplicate frame appears in more than
one partition (policy to be fixed in \tbd{O-004}). The test partition is
used only for a frozen choice.

\subsection{Baseline and candidate models}
The minimal baseline is \tbd{O-006}. Candidate backbones per modality are
\tbd{O-007}. The fine-tuning strategy is \tbd{O-008}. A small network
trained from scratch is included as a reference.

\subsection{Fusion}
For inputs with more than one modality the fusion scheme is \tbd{O-009}. A
still image has no audio, so the scheme must handle a missing modality.

\subsection{Calibration and thresholds}
The threshold is selected on validation data only. The calibration method,
if any, is \tbd{O-005}. Following~\cite{guo2017calibration}, expected
calibration error is the weighted average over bins of the absolute gap
between accuracy and mean confidence; the number of bins is fixed in advance.

\subsection{Cost measurements}
For every candidate we measure parameters, forward-pass FLOPs for the input
size actually used, peak GPU memory in a training step, seconds per epoch,
epochs to a validation target fixed in advance, sensitivity to the learning
rate, and the number of hyperparameters that had to be tuned.

\subsection{Heart-rate evidence (candidate component)}
Whether remote photoplethysmography enters the system is decided by a
feasibility test with success and failure thresholds written before the run
(\tbd{O-010}).

%%%%%%%%%%%%%%%%%%%%%%%%%%%%%%%%%%%%%%%%%%%%%%%%%%%%%%%%%%%%
\section{Experimental Protocol}\label{sec:exp}

\subsection{Setup}
Hardware: one consumer GPU with 12~GB of memory, as measured on the project
machine (environment ENV-DESKTOP-GPU; full record in the experiment log).
Software versions: \fromexp{E-xxx, environment record}.

\subsection{Comparison rules}
All candidates are compared on the same data, split, compute budget and
metric, with at least three random seeds. We report the mean, the standard
deviation over seeds, and a bootstrap confidence interval that resamples
independent units (identities or source videos), not frames. Candidates are
compared in a paired way, on the same test units. If the interval of the
difference includes zero, the simpler or cheaper candidate is preferred.
The number of resamples and the interval level are \tbd{fixed in Stage 2}.

\subsection{Metrics}
Metrics are \tbd{O-005}. Each metric is defined mathematically before use.
Class balance of every test set is stated.

\subsection{Iteration plan}
Table~\ref{tab:iters} lists the slots. Each hypothesis, success criterion
and failure criterion is written in the experiment log before its run.

\begin{table}[t]
\centering
\caption{Iteration slots (structure only; to be filled from the experiment log).}
\label{tab:iters}
\small
\begin{tabular}{@{}lp{2.2cm}p{1.6cm}p{1.6cm}@{}}
\toprule
Slot & Hypothesis & Experiment & Outcome \\
\midrule
I-0 (MVP) & \tbd{} & E-xxx & \tbd{} \\
I-1 & \tbd{} & E-xxx & \tbd{} \\
I-2 & \tbd{} & E-xxx & \tbd{} \\
I-3 & \tbd{} & E-xxx & \tbd{} \\
I-4 & \tbd{} & E-xxx & \tbd{} \\
I-5 & \tbd{} & E-xxx & \tbd{} \\
I-6 & \tbd{} & E-xxx & \tbd{} \\
I-7 & \tbd{} & E-xxx & \tbd{} \\
\bottomrule
\end{tabular}
\end{table}

%%%%%%%%%%%%%%%%%%%%%%%%%%%%%%%%%%%%%%%%%%%%%%%%%%%%%%%%%%%%
\section{Results}\label{sec:results}
% EMPTY ON PURPOSE. Every number below must be copied from a logged
% experiment and carry a marker line "% E-xxx" naming the experiment.

\subsection{Minimal baseline and iterations}
\fromexp{E-xxx, I-0 to I-7}. % E-xxx

\subsection{Final comparison}
\fromexp{E-xxx, final comparison under the common protocol}. % E-xxx

\begin{table}[t]
\centering
\caption{Final comparison on the held-out partition (structure only).
Protocol: \tbd{split name; in-dataset or cross-dataset}. Cells to be filled from the experiment IDs shown.}
\label{tab:final}
\small
\begin{tabular}{@{}lccc@{}}
\toprule
Candidate & Metric (mean $\pm$ sd) & 95\% interval & Cost \\
\midrule
\textbf{\textcolor{red}{[TBD]}} (cand.\ A) & E-xxx & E-xxx & E-xxx \\
\textbf{\textcolor{red}{[TBD]}} (cand.\ B) & E-xxx & E-xxx & E-xxx \\
\bottomrule
\end{tabular}
\end{table}

\subsection{Calibration and cross-dataset behaviour}
\fromexp{E-xxx}. % E-xxx

\subsection{Heart-rate feasibility test}
\fromexp{E-xxx, Stage 6}. % E-xxx

%%%%%%%%%%%%%%%%%%%%%%%%%%%%%%%%%%%%%%%%%%%%%%%%%%%%%%%%%%%%
\section{Discussion}\label{sec:disc}
% Written after the results. It must say what each result supports and what
% it does not, and name the experiment for every claim.
\tbd{interpretation of E-xxx results; compare differences, not boasts}.

%%%%%%%%%%%%%%%%%%%%%%%%%%%%%%%%%%%%%%%%%%%%%%%%%%%%%%%%%%%%
\section{Limitations}\label{sec:limits}
Limits that are known before any run:
\begin{itemize}
  \item Published numbers come from different protocols, so any comparison
        with them is made only by re-running a published detector on our
        protocol, not by quoting its paper.
  \item A detector trained on one kind of manipulation supports no claim
        about other kinds (for example fully generated images or cloned
        voices) without a test on them.
  \item Labelled audio-visual datasets are fewer and smaller than visual
        ones, which limits what the audio branch can support.
  \item Compute is limited to one 12~GB GPU, which bounds model size and
        clip length.
  \item \tbd{limits that appear during the experiments, from the experiment log}.
\end{itemize}

%%%%%%%%%%%%%%%%%%%%%%%%%%%%%%%%%%%%%%%%%%%%%%%%%%%%%%%%%%%%
\section{Conclusion}\label{sec:concl}
\tbd{conclusion, after results exist: what was tried, what worked, what did
not, and what the evidence does not support}.

%%%%%%%%%%%%%%%%%%%%%%%%%%%%%%%%%%%%%%%%%%%%%%%%%%%%%%%%%%%%
\bibliographystyle{IEEEtran}
\bibliography{refs}

\end{document}
